\chapter{Reading Real Greek}

Everything so far has been prepared text. In this chapter you meet
continuous, unadapted ancient Greek, read as it has come down to us for two
and a half thousand years.

\begin{figure}[ht]
  \centering
  \includegraphics[width=0.94\linewidth]{assets/greek-manuscript-6c.png}
  \caption{A sixth-century Greek cursive hand. Textual critics read
  manuscripts like this to establish the ancient text; you now know enough
  letters to begin to recognize words in it.}
\end{figure}

\section{A passage to read}

Read the following adapted sentence aloud before translating it:

\begin{quotation}
\PSEGreek{ὁ} \PSEGreek{ἄνθρωπος} \PSEGreek{ἐν} \PSEGreek{τῇ} \PSEGreek{οἰκίᾳ} \PSEGreek{γράφει} \PSEGreek{τὸν} \PSEGreek{λόγον} \PSEGreek{τῷ} \PSEGreek{βιβλίῳ}.
\end{quotation}

Even this small sentence has a story: the dative \PSEGreek{βιβλίῳ} shows the receiving
end of the writing. Every ending you learn makes the next sentence easier.

\section{Where to go next}

Keep a reader of adapted Xenophon at hand, use a dictionary such as the
\emph{Middle Liddell}, and reread your first chapters when the going gets
hard. Your steadiness will compound.\footnote{Readers of English may enjoy
the free scans of public-domain school editions, many of which are exactly
the right level after this book.}

\PSEIndex{reading practice}
\PSEIndex{manuscripts}
